\subsubsection{Transition Path Sampling}
Being able to accurately determine the height of the reaction energy barrier provides chemists and materials scientists to understand the kinetics and mechanistic details at the atomistic level for chemical reactions and phase transformations. Standard path-based methods such as the Nudged Elastic Band (NEB)\cite{jonsson1998nudged} (and its extension for for solids\cite{sheppard2012generalized}), String method,\cite{weinan2002string} and action minimisation\cite{passerone2001action} all aim to determine a well-defined reaction pathway with minimum energy barrier, with the pre-defined reactant and product. In contrast, transition path sampling (TPS)\cite{bolhuis2021transition} leads to the probabilistic view on the transitions between the initial and final state, which tries to estimate the probability distribution over different transition pathways $\mathcal{X}$ that all connect the initial and final states together through MC sampling [\cref{fig:sampling_methods}(n)]. 

Consider a standard MC moves in the configurational space that generates a canonical ensemble of configurations in a Markov chain as $\bm{R}_{i}\to\bm{R}_1\to\bm{R}_2\to\cdots\to\bm{R}_{n}\to\bm{R}_f$, from which a transition pathway can be defined as the collection of all the intermediates $\{\bm{R}_1,\ldots,\bm{R}_n\}$. Since each state in the Markov chain is linked to the next one with a finite transition  probability $p(\bm{R}_i\to\bm{R}_{i+1})$, the probability distribution over the  space of the transition pathway can therefore be written as 
\begin{equation}
    \label{eq:TPS_probability}
    p(\mathcal{X})=p(\bm{R}_0)\prod_{i=1}^{n}p(\bm{R}_{i-1}\to\bm{R}_i)/\mathcal{Z}.
\end{equation}

In the most straightforward implementation, to perform MC sampling over the distribution (\ref{eq:TPS_probability}), one can start from an existing trajectory $\mathcal{X}^{(k)}$ and randomly select one frame $\bm{R}_i^{(k)}$ from it. Then the atomic positions and/or velocities of $\bm{R}_i^{(k)}$ will be perturbed to generate a new configuration $\bm{R}_i^{(k+1)}$. From $\bm{R}_i^{(k+1)}$, one can then perform both forward (towards the final state $\bm{R}_f$ and backward (towards the initial state $\bm{R}_i$) dynamic simulations to generate a new trajectory $\mathcal{X}^{(k+1)}$, which will be accepted if both $\bm{R}_i$ and $\bm{R}_f$ can be recovered. 

There are several advantages of performing TPS over the traditional path-based methods. Firstly, it can provide a more comprehensive view on the reaction mechanism by identifying an ensemble of configurations separating the initial and final states, which are particularly useful for studying multistep reactions. Secondly, it enables extracting of a richer set of kinetic information from the underlying dynamic simulations, which also give direct access to the free energies. Finally, it enables us to compute the likelihood that a given reaction intermediate will proceed to the final state before returning to the initial state, which can be a useful sanity check to identify intermediates that are most likely to be found  under realistic thermodynamic conditions. 
